\subsection{Professional project through 2nd year in ECL and internship
research}

Coming from a research background at the ENS, through the second year of ECL
(being my first year in an engineering school) I gained an overview of the way
the industry works, with a different culture, different ways of working with
many more people than the previous experience in academic research. As I plan to
work in Research and Development as a Software Engineer I needed this point of
view from an engineering school and was very pleased to have it. I looked for an
internship that was both abroad and in a company but struggled a lot and I gave
up on looking for it in industry. When machine learning is very useful in many
topics in computer science I really wanted to do have good basis in the domain.
This internship proposal in India about generative models was the perfect match
as I also wanted to have a very new cultural experience.

\subsection{Valuable conclusions for the professional project taken from the
internship}

After having a three months experience in a forein, and culturally very
different country I really missed living in France and realized how important to
me it was not to stay too long away from my social sphere. I then realized I
didn't want to look for an internship or a job in a different country and that
the location of the job was a must for a future job. 

This internship didn't really go very far in the real life applications which is
still lacking to have a good overview of the R\&D field, but I hope I will have
this overview in the next internship.

\subsection{The internship being valuable for the professional project}

This internship was a research internship and went very deep in generative
models. I don't intent to work in neuroscience but it made me go through many
concepts that can be used as machine learning algorithms make the state of the
art in many fields, especially generative algorithms. As I intent to work in 2D
and 3D algorithms in Research and Development, it is important for me to have
good insight with academic research and to be aware of the cutting edge
algorithms in these domains.

It mostly developped my research ability as I had to read many papers,
understand them, and I tried to submit one. I had to read many of them, know the
different conferences, the grades of conferences, the requirements for each one.
I went through the hard task of having a reproducible code for a paper which is
non trivial, it requires to know how to make a very rigorous code, have people
installed the right libraries, all detailed in a well written README and as much
automatic as possible.

On the coding side I learned to code better, with a good hierarchy in my code
and good comments as it was meant to be reused. I was at ease with Python before
the internship but not that much with PyTorch and running software on a
supercomputer. Those are many habits I gained and are very valuable for the
future. I also learned good practices with GitHub, I made two pull requests and
discussed in the issues section.

\subsection{Conclusion on the professional project}

For the professional project I lack familiarity with software engineering, this
is something I must have in my next internship. I plan on making a list of
fields I would like to work in (like 3D algorithms, images, etc...) and look for
the companies in Lyon that specialize in these fields. I also plan on having
more personnal projects on my GitHub page to show real intersts and professional
code I can produce.